\chapter{Neck Pain}

\section*{Definition and Overview}
Neck pain, clinically termed cervicalgia, is a highly prevalent musculoskeletal symptom characterised by discomfort, stiffness, or pain localised to the cervical spine region, extending from the superior nuchal line to the spinous process of the first thoracic vertebra. It can manifest as an acute, self-limiting episode or as a chronic, disabling condition. Neck pain can originate from any of the cervical spine's anatomical structures, including skeletal muscles, ligaments, facet joints, intervertebral discs, and neural elements. While the vast majority of cases represent non-specific mechanical pain, neck pain can also arise from degenerative disc disease, cervical spondylosis, cervical radiculopathy, trauma (such as whiplash), or, less commonly, serious systemic pathologies.

\section*{Epidemiology}
Neck pain is a major global public health concern, consistently ranking as the fourth leading cause of years lived with disability worldwide. The lifetime prevalence of neck pain in the general population ranges between 50\% and 70\%, with an annual prevalence of 30\% to 50\%. It is more frequently reported in females than males, with a peak incidence observed in middle age (between 40 and 60 years). In Australia, epidemiological data indicate that neck pain is heavily associated with sedentary occupations, prolonged computer and smartphone use, and elevated levels of occupational stress.

\section*{Aetiology and Pathophysiology}
The underlying mechanisms of neck pain are complex and multifactorial, incorporating mechanical, degenerative, traumatic, and psychosocial elements.
\begin{itemize}
    \item \textbf{Cervical spondylosis and degenerative disc disease:} Age-related degeneration of the cervical intervertebral discs leads to decreased disc height, altered biomechanics, osteophyte formation, and subsequent stress on the facet joints.
    \item \textbf{Postural overload and muscle imbalance:} Prolonged static neck flexion (commonly referred to as "text neck") overloads the posterior cervical paraspinal muscles and upper trapezius, causing ischaemic muscle pain and myofascial trigger points.
    \item \textbf{Whiplash-associated disorders (WAD):} An acceleration-deceleration mechanism of energy transfer to the neck, typically from motor vehicle collisions, resulting in soft-tissue sprains, facet joint microtrauma, and neurogenic inflammation.
    \item \textbf{Cervical radiculopathy:} Mechanical compression or inflammatory irritation of a cervical nerve root, most commonly caused by a herniated disc or osteophyte encroaching on the intervertebral foramen.
    \item \textbf{Cervical myelopathy:} Compression of the cervical spinal cord, frequently due to advanced spinal canal stenosis or large disc herniation, leading to upper motor neurone dysfunction.
\end{itemize}

\section*{Clinical Features}
The clinical features of neck pain vary according to the anatomical structure involved and the presence of neurological compromise.
\begin{itemize}
    \item \textbf{Localised pain and stiffness:} A dull ache in the posterior neck region, often accompanied by a restricted range of motion, particularly during rotation or lateral flexion.
    \item \textbf{Referred musculoskeletal pain:} Pain radiating to the occiput, shoulder girdle, upper back, or interscapular region, which does not follow a specific nerve pathway.
    \item \textbf{Radicular pain (Cervical radiculopathy):} Sharp, shooting, or electric-shock-like pain that radiates down the arm, following a specific dermatome, often accompanied by paraesthesias, numbness, or focal muscle weakness.
    \item \textbf{Cervicogenic headache:} A unilateral, non-throbbing headache that originates in the neck and is precipitated by neck movements or sustained awkward postures.
    \item \textbf{Myelopathic features:} Clumsiness of the hands, difficulty performing fine motor tasks, gait instability, lower limb weakness, and bowel or bladder dysfunction.
\end{itemize}

\section*{Differential Diagnosis}
Clinicians must differentiate common mechanical neck pain from systemic inflammatory conditions, referred visceral pain, and serious structural diseases.
\begin{itemize}
    \item \textbf{Mechanical neck pain:} Non-specific cervical pain without radicular signs, myelopathy, or red flag symptoms.
    \item \textbf{Peripheral nerve entrapment:} Conditions such as carpal tunnel syndrome or cubital tunnel syndrome, which present with hand paresthesias but lack cervical spine tenderness or reproduction of symptoms on neck movement.
    \item \textbf{Inflammatory arthropathies:} Conditions like rheumatoid arthritis or ankylosing spondylitis, which can affect the cervical spine, causing chronic pain, morning stiffness, and atlantoaxial instability.
    \item \textbf{Vertebral artery dissection:} A rare but critical cause of acute neck pain and headache, often accompanied by neurological deficits (dizziness, double vision, dysarthria).
    \item \textbf{Spinal infection or malignancy:} Discitis, epidural abscess, or primary/metastatic spinal bone lesions presenting with persistent night pain, fever, and constitutional symptoms.
\end{itemize}

\section*{When to Seek Medical Care}
While minor neck pain is common, specific clinical red flags require immediate emergency intervention.
\textbf{Call 000 (or local emergency services) for:}
\begin{itemize}
    \item Neck pain following high-energy trauma (such as a motor vehicle crash, diving accident, or fall from a height) where cervical spine instability is suspected.
    \item Neck pain accompanied by sudden weakness, numbness, or loss of coordination in both arms or legs (quadriparesis).
    \item Sudden, severe neck pain associated with a thunderclap headache, visual disturbances, dizziness, or slurred speech.
    \item Neck pain accompanied by high fever, severe photophobia, confusion, and a stiff neck.
    \item Neck pain associated with chest pain, sweating, or shortness of breath.
\end{itemize}
Patients should consult their GP or a physiotherapist if the neck pain persists for more than 4 to 6 weeks, radiates down the arm, is accompanied by progressive numbness or tingling in the hand, or is associated with unexplained weight loss, night sweats, or a history of cancer.

\section*{Diagnosis and Investigations}
A thorough physical and neurological examination is the foundation of diagnosis. Routine imaging is not recommended for acute mechanical neck pain unless red flags are present.
\begin{itemize}
    \item \textbf{Neurological examination:} Assessing upper limb reflexes, sensory dermatomes, and myotomal muscle strength to detect nerve root or spinal cord compression.
    \item \textbf{Magnetic resonance imaging (MRI):} The imaging gold standard for evaluating soft tissues, including intervertebral discs, spinal cord compression, nerve root impingement, and infection.
    \item \textbf{Plain radiographs (X-ray):} Indicated in acute trauma (following the Canadian C-spine Rules) or to assess alignment, instability (flexion/extension views), and severe degenerative changes.
    \item \textbf{Computed tomography (CT):} Provides superior detail of bony structures, useful for evaluating complex fractures or for pre-operative planning.
    \item \textbf{Inflammatory markers (ESR and CRP):} Indicated if there is clinical suspicion of a systemic rheumatological condition or spinal infection.
\end{itemize}

\section*{Management}
The primary management of mechanical neck pain is conservative, active, and multi-modal.
\begin{itemize}
    \item \textbf{Education and active recovery:} Reassuring the patient of the benign nature of mechanical neck pain, encouraging early movement, and avoiding cervical collars, which can delay recovery.
    \item \textbf{Physical therapy and exercise:} Retraining deep cervical flexor strength, improving scapular stability, and performing range-of-motion and stretching exercises.
    \item \textbf{Pharmacological management:} Short-term use of paracetamol or oral NSAIDs to facilitate movement; neuropathic agents (e.g., amitriptyline, pregabalin) may be used for radicular pain.
    \item \textbf{Manual therapy:} Gentle joint mobilisation or soft-tissue techniques performed by qualified therapists can provide short-term symptomatic relief.
    \item \textbf{Surgical consultation:} Indicated for progressive or severe neurological deficits (such as cervical myelopathy) or intractable radicular pain that fails to respond to 6--12 weeks of conservative care.
\end{itemize}

\section*{Complications}
Chronic neck pain can lead to secondary physical and psychological complications.
\begin{itemize}
    \item \textbf{Chronic pain syndrome:} Central sensitisation resulting in persistent pain, kinesiophobia (fear of movement), and hyperalgesia.
    \item \textbf{Permanent neurological deficit:} Untreated cervical myelopathy can lead to irreversible spinal cord injury, causing spastic tetraparesis and gait impairment.
    \item \textbf{Chronic headaches:} Persistent cervicogenic headaches causing significant cognitive disruption and daily impairment.
    \item \textbf{Psychosocial distress:} Long-term neck pain is highly correlated with sleep disturbances, anxiety, depression, and reduced occupational capacity.
\end{itemize}

\section*{Prognosis and Outcomes}
The prognosis for acute mechanical neck pain is highly favourable, with the majority of patients experiencing significant recovery within 4 to 6 weeks. However, recurrence is common, with approximately 50\% of patients reporting mild recurrent episodes within a year. Whiplash-associated disorders have a more variable prognosis, with 20\% to 30\% of patients progressing to chronic pain and disability. Early adoption of active movement strategies and addressing psychosocial barriers (such as stress and fear-avoidance behaviours) are key determinants of positive long-term outcomes.

\section*{Prevention and Self-Care}
Self-management strategies can prevent the onset or recurrence of neck pain by reducing mechanical strain.
\begin{itemize}
    \item \textbf{Ergonomic workstation setup:} Positioning computer screens at eye level, using a chair with armrests to relieve shoulder strain, and positioning keyboards to allow relaxed shoulders.
    \item \textbf{Regular postural breaks:} Taking short breaks every 30--45 minutes during desk work to perform gentle neck movements and shoulder rolls.
    \item \textbf{Postural exercises:} Incorporating chin tucks and scapular retractions into daily routines to strengthen postural muscles.
    \item \textbf{Optimised sleeping posture:} Using a single, supportive pillow that aligns the neck with the rest of the spine in side or back-sleeping positions.
    \item \textbf{General physical activity:} Engaging in regular aerobic exercise (such as walking or swimming) to improve overall tissue vascularity and stress resilience.
\end{itemize}

\section*{Sources and Further Reading}
The following reputable organisations provide further information and clinical guidelines on neck pain:
\begin{itemize}
    \item Healthdirect Australia. \textit{Information on neck pain causes, symptom relief, and treatments.} https://www.healthdirect.gov.au/neck-pain
    \item Better Health Channel. \textit{Guidance on neck pain, posture, and self-care strategies.} https://www.betterhealth.vic.gov.au/
    \item NHS. \textit{Clinical guide to managing neck pain, exercises, and when to seek help.} https://www.nhs.uk/conditions/neck-pain-and-stiffness/
    \item National Institute for Health and Care Excellence (NICE). \textit{Clinical Knowledge Summaries on neck pain and cervical radiculopathy.} https://cks.nice.org.uk/
    \item Mayo Clinic. \textit{Patient resource on neck pain diagnosis, imaging, and therapeutic interventions.} https://www.mayoclinic.org/
\end{itemize}
